
\section{Superfícies de água e materiais móveis}

Uma aplicação prática do UV mapping dinâmico é a animação de superfícies que precisam transmitir sensação de movimento. Uma textura pode ser deslocada continuamente para simular o fluxo de água ou outros materiais.

A captura de uma única imagem pode mostrar a aparência da superfície, mas não demonstra completamente o movimento. Nesse caso, um vídeo de demonstração é uma evidência complementar mais adequada.

\section{Terrenos e ambientes virtuais}

Terrenos podem utilizar texturas repetíveis e procedimentos para gerar variações de aparência. Métodos procedurais também são aplicados à geração de elementos de mundos virtuais. Smelik et al. \cite{smelik2014} discutem aplicações procedurais em terrenos, vegetação, rios, estradas, edifícios e cidades.

\section{Materiais procedurais}

Materiais como pedra, madeira, metal ou superfícies com padrões repetitivos podem ser produzidos por funções procedurais. Dong et al. \cite{dong2020} revisam métodos de geração procedural de texturas bidimensionais, mostrando a relevância de modelos matemáticos para produzir diferentes aparências.

\section{Efeitos visuais}

A combinação de transformação UV com tempo permite produzir efeitos de energia, lava, brilho, movimento de padrões e outros efeitos. A vantagem é que a animação pode ser obtida modificando os parâmetros de renderização, sem necessariamente alterar a geometria do objeto.
